\section{From Saragossa to the Third Typical Edition}\label{sec:history}

Advent did not enter the Church fully formed as a four-week season with a three-year lectionary. Like Lent before it, the preparation for the Nativity grew from local vigils, episcopal fasts, and regional liturgies across Spain, Gaul, northern Italy, and Rome before converging on the Roman calendar and being renewed after the Second Vatican Council.

\subsection{Fourth-Century Spain and Fifth-Century Gaul}

The earliest surviving western conciliar record of a December season of devotion comes from fourth-century Spain. Around the year 380, twelve bishops gathered at the Council of Saragossa to counter Priscillianist tendencies toward private asceticism cut loose from the assembly. In its fourth canon, the council mandated twenty-one continuous days of daily church attendance from December~17 (\textit{a sexto decimo Kalendas Januarii}) through the Epiphany on January~6, forbidding anyone to absent himself from the church, hide at home, retreat to a country estate or the mountains, or walk barefoot in sectarian rigor. The date December~17 is striking: long before Rome composed the Great O~Antiphons, western Christians already recognized the final days before the Nativity and Epiphany as a single unbroken stretch of ecclesial gathering.

A century later in Gaul, preparation for Christmas took on an explicitly penitential and forty-day shape. St.~Gregory of Tours records in the closing chapter of his \textit{History of the Franks} (Book~X, ch.~31) that St.~Perpetuus, who governed the see of Tours around 460--490, established a written rule of fasts and vigils for the year, prescribing three fast days each week from the burial feast of St.~Martin (November~11) until the Nativity of the Lord (\textit{de depositione domni Martini usque natale Domini terna in septimana ieiunia}). Because November~11 to December~25 spans roughly forty-three days, this Gallican observance came to be known as ``St.~Martin's Lent''---a parallel to the Byzantine forty-day Nativity Fast beginning after the feast of St.~Philip on November~14, which Dom Prosper Gu\'eranger later noted as the enduring eastern counterpart of the same impulse.

Sixth-century Gallican councils confirmed and widened that observance. The Second Council of Tours (567), in its canon on monastic fasts (canon~18 in Maassen's critical numbering; canon~17 in Hefele), directed monks to fast three days a week through September, October, and November, and then every day throughout December until Christmas (\textit{de Decembre usque natale Domini omni die ieiunent}). Roughly fifteen years later, the First Council of M\^acon (581/583, canon~9) extended the St.~Martin's fast across the wider Frankish church: from the feast of St.~Martin to Christmas, Mondays, Wednesdays, and Fridays were to be kept as fast days, the Eucharistic sacrifice was to be offered after the Lenten pattern, and the canons were to be read aloud so that none could plead ignorance.

\subsection{Rome, the Tenth-Month Fast, and Five Weeks Becoming Four}

Rome arrived at Advent along a different path. In the fifth century, Pope St.~Leo the Great preached a celebrated cycle of sermons for the \textit{ieiunium decimi mensis}---the Fast of the Tenth Month (December, counting from March)---linking the winter Ember fast of Wednesday, Friday, and Saturday vigil at St.~Peter's to the gathering of the year's harvest, the defense of widows and orphans, and the discipline of desire. The earliest Roman collection of Mass prayers, the \textit{Sacramentarium Veronense}, transmits tenth-month fast formulas and Christmas prayers without yet naming a separate season called Advent.

By the seventh and eighth centuries, however, Roman sacramentaries display an organized Advent proper, though with a notable variation in length. The Old Gelasian Sacramentary (\textit{Sacramentarium Gelasianum Vetus}, Book~II, sections LXXX--LXXXIV) provides five Advent Sunday Masses under the heading \textit{Orationes de Adventum Domini}, opening with the stirring cry \textit{Excita, Domine, potentiam tuam, et veni} (``Stir up, O Lord, thy power, and come''), and praying in the second Mass that God may purify our consciences by visiting them daily so that His Son at His coming may find a dwelling prepared within us. Five Advent Sundays likewise appear in the ancient Spanish \textit{Liber Mozarabicus Sacramentorum} (later expanded to six in printed Mozarabic books), while Milan's Ambrosian Rite preserved a six-Sunday Advent beginning after St.~Martin's feast (\textit{Manuale Ambrosianum}, six Sundays \textit{de Adventu}).

Rome itself settled on four Sundays, as witnessed in the Gregorian Sacramentary sent by Pope Hadrian~I to Charlemagne (\textit{Sacramentarium Gregorianum Hadrianum}). There, under \textit{Mense Decembrio Orationes de Aduentu Domini}, appear \textit{Dominica prima}, \textit{Dominica secunda}, \textit{Dominica III ad sanctum Petrum}, and then the famous rubric \textit{Die Dominica uacat} (``the Sunday is vacant'')---because the night-long Ember Saturday ordination vigil at St.~Peter's ran into Sunday morning---followed by the \textit{Excita} collect that later anchored the Fourth Sunday of Advent. In the ninth century, the Carolingian liturgist Amalarius of Metz (\textit{Liber Officialis} III.40; \textit{Liber de Ordine Antiphonarii} XXX) already remarked on the shift from five-week lectionaries to the four-week Roman order attributed to Pope St.~Gregory the Great, and noted that from ancient times the \textit{Gloria in excelsis Deo} was omitted during Advent to signify the watchful expectation of the Old Testament.

Alongside the Mass prayers, Rome and the Frankish monasteries wove the Advent Office from Isaiah and the Gospels, crowning the final week before Christmas Eve with the seven Great O~Antiphons---already attested in ninth-century Carolingian liturgy and preserved in medieval chant collections such as the Hartker Antiphonary of St.~Gall. Through the medieval Roman Curia and the Tridentine \textit{Missale Romanum} (1570), this four-week Roman Advent endured down to the 1962 typical edition with its four Sunday Mass sets, its Ember Wednesday, Friday, and Saturday Masses, and its Vigil of the Nativity.

\subsection{The Second Vatican Council and the Postconciliar Renewal}

The Second Vatican Council's Constitution on the Sacred Liturgy, \textit{Sacrosanctum Concilium} (4 December 1963), set the principles from which the postconciliar Advent was shaped. In nos.~5--10, the Council grounded all liturgy in the paschal mystery of Christ, the summit toward which the Church's activity is directed and the font from which all her power flows. In no.~50, it called for the revision of the Order of Mass so that the proper nature and interconnection of each part would shine clearly and invite the full, conscious, and active participation of the faithful (no.~14). In no.~51, it mandated that the treasures of the Bible be opened more lavishly so that a richer fare of God's word would be spread before the people over a fixed cycle of years. In no.~91, it asked that the Psalter of the Divine Office be distributed over a longer period than a single week. And in nos.~102--111, it ordered a revision of the liturgical year so that the traditional customs and discipline of the sacred seasons would be preserved or restored to suit the conditions of modern times, keeping Sunday and the temporal cycle supreme over lesser feasts of saints (no.~108).

Those conciliar mandates took concrete canonical and textual form in four landmark promulgations under Pope St.~Paul~VI:

\begin{itemize}
  \item \textbf{The Revised Calendar (1969).} On 14 February 1969, by the motu proprio \textit{Mysterii Paschalis}, Pope Paul~VI approved the \textit{Universal Norms on the Liturgical Year and the Calendar} and the new General Roman Calendar (effective 1 January 1970). Nos.~39--42 gave Advent its explicit two-phase charter, while nos.~45--47 entrusted Ember and Rogation days to the pastoral judgment of bishops' conferences.
  \item \textbf{The Revised Missal (1969/1970).} On 3 April 1969, by the apostolic constitution \textit{Missale Romanum}, Paul~VI promulgated the new Roman Missal; the Sacred Congregation for Divine Worship issued the first typical edition on 26 March 1970, followed by the second typical edition in 1975. The Missal supplied proper orations for every single weekday of Advent, paired with two proper Advent Prefaces (Preface~I for the first phase, looking to the glorious return; Preface~II for December~17--24, looking to the imminent Nativity) and expanded weekday antiphons drawn from ancient Roman and Breviary chants.
  \item \textbf{The Revised Lectionary (1969/1981).} Promulgated in 1969 and issued in its expanded second typical edition (\textit{Ordo Lectionum Missae}, 1981), the Lectionary created the three-year Sunday cycle (Years~A, B, and~C) and the two-part weekday Advent course: seventeen week-numbered weekdays of Isaiah and the Gospels for Weeks~I--III (nos.~175--191, followed by the Advent weekday Gospel Acclamation verses at no.~192) and eight date-numbered privileged weekdays for December~17--24 (nos.~193--200).
  \item \textbf{The Liturgy of the Hours (1970).} On 1 November 1970, by the apostolic constitution \textit{Laudis canticum}, Paul~VI promulgated \textit{Liturgia Horarum}, reshaping Morning and Evening Prayer for clergy, religious, and laity alike, introducing a four-week Psalter, and assigning Isaiah and rich patristic treatises to the Advent Office of Readings.
\end{itemize}

\subsection{The Third Typical Edition and the 2011 United States Missal}

During the Great Jubilee of 2000, Pope St.~John Paul~II approved the third typical edition of the \textit{Missale Romanum} (decree of the Congregation for Divine Worship and the Discipline of the Sacraments, 20 April 2000; published in 2002), incorporating decades of calendar updates---including universal memorials such as St.~Juan Diego (December~9) and Our Lady of Guadalupe (December~12, elevated to a Feast in the United States)---together with a revised \textit{General Instruction of the Roman Missal} (GIRM) and expanded blessings. In 2008, the Holy See issued the \textit{reimpressio emendata} with targeted textual corrections published in \textit{Notitiae}, and in 2019 Pope Francis inscribed the optional memorial of the Blessed Virgin Mary of Loreto on December~10.

In the United States, following the translation norms of \textit{Liturgiam authenticam} (2001), the English translation of the Third Typical Edition of \textit{The Roman Missal} entered into liturgical use on the First Sunday of Advent, 27 November 2011, paired with the United States \textit{Lectionary for Mass} (second typical edition, 1998--2002) and the four-volume \textit{Liturgy of the Hours} (1975). Across more than sixteen centuries---from Saragossa's twenty-one December days and St.~Leo's tenth-month vigils to the third typical edition on the parish altar---the Church's Advent has matured without ever losing its center: the watch of a pilgrim people waiting for the Lord who came in humility, comes now in grace, and will come again in glory.
